
\subsection{Proof of Lemma~\ref{lem:abscfg_sound}}
\label{apdx:abscfg_sound}
\begin{lemma}[Soundness of the Abstract Events Computation]
  For any program $c$ with 
  % its abstract control flow graph $\absG(c) = (\absV(c), \absE(c))$,
  an arbitrary initial trace $\vtrace_0 \in \ftdom_0(c)$ and
  a finite execution trace $\trace \in \ftdom$ that is generated 
  when $c$ is executing under $\vtrace_0$,
  then for every event $\event \in \trace$ such that $\event = (\_, l, \_, \_)$ for some label $l$ in program $c$,
  % an initial trace  $\vtrace_0 \in \ftdom_0(c)$,
  there is an abstract event $\absevent = (l, \_, \_) \in \absflow(c)$ and $\absevent = (l, \_, \_)$.
%
\[
\begin{array}{l}
  \forall c, c' \in \cdom, \trace_0 \in \ftdom_0(c), \trace \in \ftdom ,  \event = (\_, l, \_, \_) \in \eventset, l \in \progl(c) \sthat
  \\
  \qquad
    \big(
    \config{{c}, \trace_0} \to^{*} \config{c', \trace_0 \tracecat \trace} 
  % \lor 
  % \config{{c}, \trace_0} \uparrow^{\infty} \trace_0 \tracecat \trace 
  \big)
  \land \event \in \trace 
   \implies \exists 
  %  \in (\ldom\times \dcdom^{\top} \times \ldom) \sthat 
  \absevent \in \absflow(c) \land \absevent = (l, \_, \_)
\end{array}
\]
\end{lemma}

\begin{proof}
  Taking an arbitrary program $c$, an intermediate program $c' \in \cdom$, an initial trace $\trace_0 \in \ftdom_0(c)$,  a finite execution trace $\trace \in \ftdom$, a label $l \in \progl(c)$ and an event $\event = (\_, l, \_, \_) \in \eventset$
  such that
  \[
    \config{{c}, \trace_0} \to^{*} \config{c', \trace_0 \tracecat \trace}
    % \lor 
    % \config{{c}, \trace_0} \to^{\infty} \config{c', \trace_0 \tracecat \trace}.
  \]
  Then it is sufficient to show,
  \[
    \event \in \trace
    \implies
    \exists \absevent \in \absflow(c) \land \absevent = (l, \_, \_)
  \]
  By induction on program $c$, we have the following cases:
  \caseL{$c = \clabel{\assign{x}{\expr}}^{l'}$}
  By the evaluation rule $\rname{assn}$, we have
  \[
  {
  \config{\clabel{\assign{{x}}{\aexpr}}^{l'},  \trace_0 } 
  \xrightarrow{} 
  \config{\clabel{\eskip}^{l'}, \trace_0 :: ({x}, l', v, \bullet)}
  }\], for some $v \in \mathbb{N}$.
  \\
  There are 2 cases, where $l' = l$ and $l' \neq l$.
  \\
  In case of $l' \neq l$, we know $\event \not\in [({x}, l', v, \bullet)]$, then this Lemma is vacuously true.
  \\
  In case of $l' = l$, by the abstract execution trace computation in Definition~\ref{def:absevent_compute}, we have 
  \[
    \absflow(c) = \absflow'(\clabel{\assign{{x}}{\aexpr}}^{l'}; \clabel{\eskip}^{\lex}) = \{(l', \absexpr(\expr), \lex)\}
  \]
  Let $\absevent = (l', \absexpr(\expr), \lex) $, then we have $\absevent \in \absflow(c)$ and this case is proved.
  \caseL{$c = \clabel{\assign{x}{\query(\qexpr)}}^{l'}$}
  By the evaluation rule $\rname{query}$
      \[
          \inferrule
          {
            \config{\trace, \qexpr} \qarrow \qval
           \and 
          \query(\qval) = v
          \and 
          \event = ({x}, l', v, \qval)
          }
          {
          \config{{\clabel{\assign{x}{\query(\qexpr)}}^{l'}, \trace}}
          \xrightarrow{} 
          \config{{\clabel{\eskip}^{l'},  \trace \traceadd \event} }
          },
        \]
      let $v \in \mathbb{N}$ be the result of query request and $\qval$ be the query value.
      Then we have
      \[
      {
        \config{{\clabel{\assign{x}{\query(\qexpr)}}^{l'}, \trace_0}}
        \xrightarrow{} 
        \config{{\clabel{\eskip}^{l'},  \trace_0 \traceadd [({x}, l', v, \qval)]} }.
      }\]
  \\
  There are 2 cases, where $l' = l$ and $l' \neq l$.
  \\
  In case of $l' \neq l$, we know $\event \not\in [({x}, l', v, \qval)]$, then this Lemma is vacuously true.
  \\
  In case of $l' = l$, by the abstract execution trace computation in Definition~\ref{def:absevent_compute}, we have 
  \[
    \absflow(c) = \absflow'(\clabel{\assign{x}{\query(\qexpr)}}^{l'}; \clabel{\eskip}^{\lex}) = \{(l', \absexpr(\expr), \lex)\}
  \]
  Let $\absevent = (l', \absexpr(\expr), \lex) $, then we have $\absevent \in \absflow(c)$ and this case is proved.
   \caseL{$\ewhile \clabel{b}^{l_w} \edo (c_w)$}
  By the first rule applied to $c$ in operational semantics, there are two cases as follows.
  \subcaseL{$\rname{while-t}$}
  In this case, we have
  \[
    \config{{\ewhile \clabel{b}^{l_w} \edo (c_w), \trace_0}}
    \xrightarrow{} 
    \config{{
    c_w; \ewhile \clabel{b}^{l_w} \edo (c_w),
    \trace_0 :: (b, l_w, \etrue, \bullet)}}.
  \]
  In case of $l \neq l_w$, we know $\event \not \in [(b, l_w, \etrue, \bullet)]$, and this Lemma is vacuously true.
  \\
  In case of $l = l_w$, let $\event = (b, l_w, \etrue, \bullet)$ and $\trace' =  [(b, l_w, \etrue, \bullet)]$, then it is sufficient to show.
    \[
      \exists 
      \absevent \in \absflow(c) \land \absevent = (l, \_, \_)
    \]
    By abstract trace computation in Definition~\ref{def:absevent_compute}, we have
    \[
      \absflow(c) = \absflow'(c; \clabel{\eskip}^{\lex}) = \{(l_w, \top, \init(c_w))\} \cup \cdots.
    \]
    Let $\absevent = (l_w, \top, \init(c_w))$, then we have this case is proved.
  \subcaseL{$\rname{while-f}$}
  In this case, we have
  \[
    \config{{\ewhile \clabel{b}^{l_w} \edo (c_w), \trace_0}}
    \xrightarrow{}^\rname{while-f}
    \config{{
      \clabel{\eskip}^{l_w},
    \trace_0 :: (b, l_w, \efalse, \bullet)}}.
  \]
  In case of $l \neq l_w$, we know $\event \not\in [(b, l_w, \efalse, \bullet)]$, and this Lemma is vacuously true.
  \\
  In the case where $l = l_w$, 
  let $\trace = [(b, l_w, \efalse, \bullet)]$ and $\event = (b, l_w, \efalse, \bullet)$,
  then it is sufficient to show
    \[
      \event \in [(b, l_w, \efalse, \bullet)]
      \implies
      \exists \absevent = (l, \_, \_) \in (\ldom\times \dcdom^{\top} \times \ldom) \sthat 
      \absevent \in \absflow(c).
    \]
  By abstract trace computation in Definition~\ref{def:absevent_compute}, we have
  \[
    \absflow(c) = \absflow'(c; \clabel{\eskip}^{\lex}) = \{(l_w, \top, \init(c_w))\} \cup \cdots.
  \]
  Let $\absevent = (l_w, \top, \init(c_w))$, then we have this case is proved.
  \caseL{$\eif(\clabel{b}^{l}, c_t, c_f)$}
  By the first rule applied to $\eif(\clabel{b}^l, c_t, c_f)$ in operational semantics,
  we have two subcases \rname{if-t} and \rname{{if-f}}.
  \\
  Then both of the two cases will append the evaluation result of the guard $\clabel{b}^l$, i.e., $(b, l, \etrue, \bullet)$ and $(b, l, \efalse, \bullet)$ respectively.
  \\
  By abstract trace computation in Definition~\ref{def:absevent_compute}, we have
  \[
    \begin{array}{l}
      \absflow(c) = \absflow'(c; \clabel{\eskip}^{\lex}) 
      \\ \quad
      = \absflow'(\eif \clabel{b}^l \ethen c_t \eelse c_f) \cup \absflow'(\clabel{\eskip}^{\lex}) 
      \\ \quad \quad 
      \cup \{ (l, dc, \absinit(c_2)) | (l, dc) \in \absfinal(c_1) \} 
      \\ \quad  
      = \{(l, \absexpr(\bexpr, \top),  \absinit(c_t) ) ,  (l, \absexpr(\neg\bexpr, \top), \absinit(c_f)) \}
      \\ \quad \quad 
      \cup \absflow'(c_t) \cup \absflow'(c_f) 
      \cup \{ (l, dc, \lex) | (l, dc) \in \absfinal(c_1) \} 
   \end{array}
  \]
  Let $\absevent = (l, \absexpr(\bexpr, \top),  \absinit(c_t) ) $ or $ (l, \absexpr(\neg\bexpr, \top), \absinit(c_f)) $
  respectively in the two subcases.
  \\
  Then we have both of the two subcases proved.
  \caseL{$c = c_1;c_2$}
  By the first rule applied to $c$ in operational semantics, there are two cases as follows.
  \subcaseL{$\rname{seq1}$}
  In this case, we have the following where $\trace = \trace_1$.
  \[
    \inferrule
    {
    \config{{c_1, \trace_0}}
    \xrightarrow{}
    \config{{c_1',  \trace_0 \tracecat \trace_1}}
    }
    {
    \config{{c_1; c_2, \trace_0}} 
    \xrightarrow{} 
    \config{{c_1'; c_2, \trace_0 \tracecat \trace_1}}
    }
  \]
  By induction hypothesis on $c_1$, we know 
  \[
    \event \in \trace_1
    \implies
    \exists 
    % \in (\ldom \times \dcdom^{\top} \times \ldom) \sthat 
    \absevent \in \absflow(c_1) \land \absevent = (l, \_, \_) 
  \]
  Let $\event = (\_, l_1, \_, \_) \in \trace_1$ and $\absevent_1 = (l_1, \_, \_) \in \absflow(c_1)$, we know $\event \in \trace$ as well.
  Then it is sufficient to show
  \[
    \exists 
    \absevent \in \absflow(c) \land \absevent = (l_1, \_, \_)
  \]
  By abstract trace computation in Definition~\ref{def:absevent_compute}, we have
  \[
    \begin{array}{l}
      \absflow(c_1) = \absflow'(c_1; \clabel{\eskip}^{\lex}) 
      \\ \quad
      = \absflow'(c_1) \cup \absflow'(\clabel{\eskip}^{\lex}) 
      \\ \qquad
      \cup \{ (l, dc, \absinit(\clabel{\eskip}^{\lex}))) | (l, dc) \in \absfinal(c_1) \} 
      \\ \quad
      = \absflow'(c_1) \cup  \{ (l, dc, \lex) | (l, dc) \in \absfinal(c_1) \} 
   \end{array}
  \]  
  and 
  \[
    \begin{array}{l}
      \absflow(c) = \absflow'(c; \clabel{\eskip}^{\lex}) 
      \\ \quad
      = \absflow'(c_1;c_2) \cup \absflow'(\clabel{\eskip}^{\lex}) 
      \\ \qquad  
      \cup \{ (l, dc, \absinit(\clabel{\eskip}^{\lex})) | (l, dc) \in \absfinal(c_1; c_2) \} 
      \\ \quad  
      = \absflow'(c_1) \cup \absflow'(c_2) 
      \\ \quad \quad
      \cup \{ (l, dc, \absinit(c_2)) | (l, dc) \in \absfinal(c_1) \} 
      \\ \quad \quad
      \cup \{ (l, dc, \lex) | (l, dc) \in \absfinal(c_1; c_2) \} 
   \end{array}
  \]
  Then we have two subcases where 
  $\absevent_1 \in \absflow'(c_1)$ or $\absevent_1 \in \{ (l, dc, \lex) | (l, dc) \in \absfinal(c_1) \}$.
  \subcaseL{$\absevent_1 \in \absflow'(c_1)$}
  Since $\absflow'(c_1) \subseteq \absflow(c)$, then we know $\absevent_1 \in \absflow(c)$.
  \\
  Let $\absevent = \absevent_1$, then we have $\absevent \in \absflow(c)$ and this subcase is proved.
  \subcaseL{$\absevent_1 \in \{ (l, dc, \lex) | (l, dc) \in \absfinal(c_1) \} $}
  In this case, we know $(l_1, \_) \in \absfinal(c_1)$.
  \\
  Let $\absevent = (l_1, \_, \absinit(c_2))$, then we know 
  \[
    \absevent \in   \{ (l, dc, \absinit(c_2)) | (l, dc) \in \absfinal(c_1) \}. 
  \]
  Since $\{ (l, dc, \absinit(c_2)) | (l, dc) \in \absfinal(c_1) \}  \subset\absflow(c)$,  then we have $\absevent \in \absflow(c)$ and this subcase is proved.
  \subcaseL{$\rname{seq2}$}
  In this case, we have the following where $\trace = \trace_2$.
  \[
    \inferrule
    {
    \config{{c_2, \trace_0}}
    \xrightarrow{}
    \config{{c_2',  \trace_0 \tracecat \trace_2}}
    }
    {
    \config{{\clabel{\eskip}^l; c_2, \trace_0}} 
    \xrightarrow{} 
    \config{{c_2, \trace_0 \tracecat \trace_2}}
    }
  \]
  By induction hypothesis on $c_2$, we know 
  \[
    \event \in \trace_2
    \implies
    \exists 
    \absevent \in \absflow(c_2) \land \absevent = (l, \_, \_) 
  \]
  Let $\event = (\_, l_2, \_, \_) \in \trace_2$ and $\absevent_2 = (l_2, \_, \_) \in \absflow(c_2)$, we know $\event \in \trace$ as well.
  Then it is sufficient to show
  \[
    \exists 
    % \in (\ldom \times \dcdom^{\top} \times \ldom) \sthat 
    \absevent \in \absflow(c) \land  \absevent = (l_2, \_, \_)
  \]
  By abstract trace computation in Definition~\ref{def:absevent_compute}, we have
  \[
    \begin{array}{l}
      \absflow(c_2) = \absflow'(c_2; \clabel{\eskip}^{\lex}) 
      \\ \quad
      = \absflow'(c_2) \cup \absflow'(\clabel{\eskip}^{\lex}) 
      \\ \qquad
      \cup \{ (l, dc, \absinit(\clabel{\eskip}^{\lex}))) | (l, dc) \in \absfinal(c_2) \} 
      \\ \quad
      = \absflow'(c_2) \cup  \{ (l, dc, \lex) | (l, dc) \in \absfinal(c_2) \} 
   \end{array}
  \]  
  and 
  \[
    \begin{array}{l}
      \absflow(c) = \absflow'(c; \clabel{\eskip}^{\lex}) 
      \\ \quad
      = \absflow'(c_1;c_2) \cup \absflow'(\clabel{\eskip}^{\lex}) 
      \\ \qquad
      \cup \{ (l, dc, \absinit(\clabel{\eskip}^{\lex})) | (l, dc) \in \absfinal(c_1; c_2) \} 
      \\ \quad  
      = \absflow'(c_1) \cup \absflow'(c_2) 
      \\ \quad \quad
      \cup \{ (l, dc, \absinit(c_2)) | (l, dc) \in \absfinal(c_1) \} 
      \\ \quad \quad
      \cup \{ (l, dc, \lex) | (l, dc) \in \absfinal(c_1; c_2) \} 
   \end{array}
  \]
  Then we have two subcases where 
  $\absevent_2 \in \absflow'(c_2)$ or $\absevent_2 \in \{ (l, dc, \lex) | (l, dc) \in \absfinal(c_2) \}$.
  \subcaseL{$\absevent_2 \in \absflow'(c_2)$}
  Since $\absflow'(c_2) \subseteq \absflow(c)$, then we know $\absevent_2 \in \absflow(c)$.
  \\
  Let $\absevent = \absevent_2$, then we have $\absevent \in \absflow(c)$ and this subcase is proved.
  \subcaseL{$\absevent_2 \in \{ (l, dc, \lex) | (l, dc) \in \absfinal(c_2) \} $}
  In this case, we know $(l_2, \_) \in \absfinal(c_2)$. By abstract continuation state computation in Section~\ref{sec:abs_prog-edge}, we know
  \[
    \absfinal(c_1; c_2) =  \absfinal(c_2) 
  \]
  Then we know $(l_2, \_) \in \absfinal(c_1; c_2)$. Let $\absevent = \absevent_2$, then we know 
  \[
    \absevent = \absevent_2 \in   \{ (l, dc, \lex) | (l, dc) \in \absfinal(c_1; c_2) \}  \subset\absflow(c).
  \]
  Then we have $\absevent \in \absflow(c)$ and this subcase is proved.
  \end{proof}


\subsection{Proof of Lemma~\ref{lem:abscfg_uniquex}}
\label{apdx:abscfg_uniquex}
\begin{lemma}[Uniqueness of the Abstract Events Computation]
  For every program $c$ and
  an execution trace $\trace \in \ftdom$ that is generated w.r.t.
  an initial trace  $\vtrace_0 \in \ftdom_0(c)$,
  there is a unique abstract event $\absevent = (l, \_, \_) \in \absflow(c)$ 
  for every assignment event $\event \in \eventset^{\asn}$ in the
  execution trace having the label $l$, i.e., $\event = (\_, l, \_, \_)$ and  $\event \in \trace$.
%
\[
  \begin{array}{l}
    \forall c \in \cdom, \vtrace_0 \in \ftdom_0(c), \trace \in \ftdom ,  l, l' \in \absV(c), \event = (\_, l, \_, \_) \in \eventset^{\asn} \sthat 
    \\ \qquad
    \config{{c}, \trace_0} \to^{*} \config{\clabel{\eskip}^{l'}, \trace_0 \tracecat \vtrace} 
  \land \event \in \trace 
    \\
    \qquad \implies \exists! \absevent = (l, \_, \_) \in (\ldom\times \dcdom^{\top} \times \ldom) \sthat  
    \absevent \in \absflow(c)
\end{array}
\]
\end{lemma}
\begin{proof}
  This is proved trivially by induction on the program $c$.
\end{proof}
%
%

\subsection{Proof of Theorem~\ref{thm:transition_bound_sound}}
  \label{apdx:weight_soundness}
  Preliminary Theorem from paper~\cite{SinnZV17}.
  \begin{theorem}[Soundness of the Transition Bound]
    % \label{thm:transition_bound_sound}
    For each program ${c}$ and an edge $\absevent = (l, \_, \_) \in \absG(c)$, if $l$ is the label of an assignment command,
    %  label $l \in \lvar(c)$,
    then its \emph{path-insensitive transition bound} $\absclr(\absevent, c)$ 
     is a sound upper bound on 
    %  execution-based reachability bound $w^t$ 
    the execution times of this assignment command in $c$.
      \[
        \begin{array}{l}
          \forall c \in \cdom, l, l' \in c,\trace_0 \in \ftdom_0(c), 
          \trace \in \tdom, v \in \mathbb{N}
           \sthat 
           \\ \qquad 
           \Big(\config{{c}, \trace_0} \to^{*} \config{\clabel{\eskip}^{l'}, \trace_0\tracecat\vtrace} 
           \land \config{\absclr(\absevent, c), \trace_0} \aarrow v
           \land
          \vcounter(\trace, l) \leq v
          \Big)
          % \\ \land \quad 
          % \Big(\config{{c}, \trace_0} \uparrow^{\infty} \trace_0\tracecat\vtrace
          %  \implies \absclr(\absevent, c) = \infty 
          %  \Big)
        \end{array}
        \]
  \end{theorem}

  {
  \begin{theorem}[Soundness of the Weight Estimation]
    % \label{thm:rb_soundness}
    Let  ${c}$ be a program and $\progW(c)$ be its estimated weight set.
    Then, for every $(x^l, \hat{w}) \in \progW(c)$ and $(x^l, w) \in \traceW(c) $,
    we have for every possible 
    $\vtrace_0 \in \ftdom_0(c),
    v \in \mathbb{N}$,
    if  $\config{\vtrace_0, \hat{w}} \earrow v$,
    then $v$ is an upper bound on ${w}(\trace_0)$.
    \[
      \begin{array}{l}
        \forall c \in \cdom, x^l \in \lvar(c),\trace_0 \in \ftdom_0(c), \trace \in \ftdom,
    v \in \mathbb{N}, l', (x^l, w) \in \traceW(c) \sthat
    \\ \qquad 
    \exists (x^l, \hat{w}) \in \progW(c) 
    \land
    % \config{{c}, \trace_0} \to^{*} \config{\clabel{\eskip}^{l'}, \trace_0 \tracecat\vtrace} 
    % \land 
    \big(
      \config{\vtrace_0, \hat{w}} \earrow v \implies w(\trace_0) \leq v
    \big)
    \end{array}
    \]
  \end{theorem}
  }
\begin{proof}
  Taking an arbitrary a program ${c}$ with its program-based dependency graph $\progG = (\vertxs, \edges, \weights, \qflag)$, 
  and an arbitrary pair of labeled variable and weights $(x^l, w) \in \progW$, 
  and arbitrary $\vtrace, \trace' \in \mathbb{T},
  v \in \mathbb{N}$ satisfying
  \\
  $\config{{c}, \trace} \to^{*} \config{\clabel{\eskip}^{l'}, \trace\tracecat\vtrace'} 
  \land 
  \config{\trace, w} \earrow v$

  
  By Definition of $\progW$ in $\progG(c)$, we know 
  $  w = \absW(l) = \max \{ \absclr(\absevent) | \absevent = (l, \_, \_)\}$.

  
  By Lemma~\ref{lem:abscfg_sound}, there exists an abstract event in $\absflow(c)$ of form $(\absevent) = (l, \_, \_)$,
  corresponding to the assignment command associated to labeled variable $x^l$. 

  
  Let $(\absevent) = (l, dc, l') \in \absflow(c)$ be this event for some $dc$ and $l'$ such that  $(\absevent) = (l, dc, l') \in \absflow(c)$,
  by the last step of phase 2, we know
  $
  \progW(x^l) 
  \triangleq \absclr(\absevent)
  $.
   Then, it is sufficient to show:
  \[
  \forall v \in \mathbb{N} \sthat  
  \config{\absclr(\absevent), \trace} \earrow 
  \vcounter(\vtrace', l) \leq v
  \absclr(\absevent)
  \]
  % By line:2 of Algorithm~\ref{alg:add_weights}, there are 2 cases:
  By definition of $\absclr(\absevent)$:
  \[
 \begin{array}{ll}
  \locbound(\absevent) & \locbound(\absevent) \in \constdom \\
  Incr(\locbound(\absevent)) + 
  \sum\{\absclr(\absevent') \times \max(\varinvar(a) + c, 0) | (\absevent', a, c) \in \reset(\locbound(\absevent))\} 
  & \locbound(\absevent) \notin \constdom
\end{array}
\]
  \caseL{$\locbound(\absevent) \in \constdom$}
  \\
  Proved by the soundness of Local bound in Lemma~\ref{lem:local_bound_sound}.
  \caseL{$\locbound(\absevent) \notin \constdom$}
To show:
\[
  \begin{array}{l}
    \max \left\{ \vcounter(\vtrace') l ~ \middle\vert~
\forall \vtrace \in \ftdom \sthat  \config{{c}, \trace} \to^{*} \config{\clabel{\eskip}^{l'}, \trace\tracecat\vtrace'} \right\} 
\\
\leq 
Incr(\locbound(\absevent)) + 
\sum\{\absclr(\absevent') \times \max(\varinvar(a) + c, 0) | (\absevent', a, c) \in \reset(\locbound(\absevent))\} 
\end{array}
\]
  Taking an arbitrary initial trace
  $\trace_0 \in \ftdom$, 
  executing $c$ with $\trace_0$, let $\trace$ be the trace after evaluation, i.e., $\config{{c}, \trace_0} \to^{*} \config{\clabel{\eskip}^{l'},\vtrace}$, it is sufficient to show:
  \[ 
    \begin{array}{l}
      \vcounter(\vtrace') l \leq 
    Incr(\locbound(\absevent)) + 
    \sum\{\absclr(\absevent') \times \max(\varinvar(a) + c, 0) | (\absevent', a, c) \in \reset(\locbound(\absevent))\}
  \end{array}
  \]
%
 By the soundness of the (1) Transition Bound and (2) Variable Bound Invariant 
 in Theorem 1 by \cite{SinnZV17} (attached above in Theorem~\ref{thm:transition_bound_sound}), 
This case is proved.
\end{proof}

\begin{lem}[Soundness of the Local Bound]
  \label{lem:local_bound_sound}
Given a program ${c}$, we have:
%
\[
\forall \absevent = (l, dc, l') \sthat  
\max \left\{ \vcounter(\vtrace') l ~ \middle\vert~
\forall \vtrace \in \ftdom \sthat  \config{{c}, \trace} \to^{*} \config{\clabel{\eskip}^{l'}, \trace\tracecat\vtrace'} \right\} 
\leq 
\locbound(\absevent)
\]
\end{lem}
\begin{proof}
  \subcaseL{$l \notin SCC(\absG(c))$}
  In this case, we know variable $x^l$ isn't involved in the body of any $\ewhile$ command. 
  \\
  Taking an arbitrary $\vtrace_0 \in \ftdom$, 
  let $\trace \in \ftdom$ be of resulting trace of executing $c$ with $\trace$, 
  i.e., $\config{{c}, \trace_0} \to^{*} \config{\clabel{\eskip}^{l'}, \trace}$,
  \\
  we know the
  assignment command at line $l$ associated with the abstract event $\absevent$ will be executed at most once, i.e.,:
  %
  $\vcounter(\vtrace) l \leq 1$
  \\
  By $\locbound$ definition, we know $\locbound(\absevent) = 1$.
  \\
  This case is proved.
  \subcaseL{$l \in SCC(\absG(c)) \land \absevent \in \dec(x) $}  in this case, we know $\locbound(\absevent) \triangleq x$.
  \subcaseL{$l \in SCC(\absG(c)) \land \absevent 
  \notin \bigcup_{x \in VAR} \dec(x)
  \land \absevent \notin SCC(\absG(c)/\dec(x)) $}  in this case, we know $\locbound(\absevent) \triangleq x$.
  \\
  In the two cases above, the soundness is discussed by \cite{SinnZV17} in the Paragraph \emph{Discussion on Soundness} in Section 4, Page 25.
\end{proof}